% SPDX-License-Identifier: Apache-2.0
\section{Several Invocations in Flight}
\label{sec:x-inflight}

The embedding article claims that a pipeline is an \code{async fn}
driven by starting an invocation every cycle, and this is the claim
run. \code{mac} has two awaits, a multiply and an add, so it takes two
cycles; \code{pipeline::drive} starts one invocation at every edge at
which a pair is offered, polls every invocation in flight as a process
of its own, and sends each result as its invocation completes, oldest
first. Each invocation prints where it is, and the trace shows what a
pipeline is: at every edge from the second, one invocation is in its
multiplier and another in its adder, and from the third a result lands
every cycle. Nobody placed a stage boundary; the two awaits are the
boundaries. Figure~\ref{fig:x-inflight} shows the pairs going in and
the results coming out two edges later.

\begin{figure}[htbp]
\centering
\input{inflight_timing.tex}
\caption{The pipeline: a pair in every cycle, a result out every cycle from the third, two invocations in flight.}
\label{fig:x-inflight}
\end{figure}

\lstinputlisting[caption={\code{ex\_inflight.rs}. Run with \code{bazel run //lib/examples:ex\_inflight}.},label={lst:x-inflight}]{lib/examples/ex_inflight.rs}

\lstinputlisting[style=ascii,caption={What \code{ex\_inflight} printed when the build ran it.},label={lst:x-inflight-out}]{out_inflight.txt}
